% GENERATED FILE. Edit canonical lessons, then run npm run generate:books.
% canonical-source-hash: fnv1a64:bf9f1f4c43843908
% canonical-lessons: GU-C250-kampyutar, GU-C250-skrin, GU-C250-dabavvu, GU-C250-sacavvu, GU-C250-chapvu

\chapter{At the computer}
\label{ch:gu-at-the-computer}
\hlchaptermodality{250}

\begin{chapteropening}
\textbf{By the end of this chapter:} \emph{I can talk about using a computer: pressing keys, saving and printing.}

Complete the last lesson of chapter 250: I can talk about using a computer: pressing keys, saving and printing.
\end{chapteropening}

\section[\texorpdfstring{\gu{કમ્પ્યૂટર} (kampyūṭar)}{કમ્પ્યૂટર (kampyūṭar)}]{\gu{કમ્પ્યૂટર} (kampyūṭar) — computer}

\label{lesson:GU-C250-kampyutar}



\begin{quote}
Before the new one: say the Gujarati for voice, then the Gujarati for to obtain.
\end{quote}

\subsection*{You'll want to know: \gu{કમ્પ્યૂટર}}
\textbf{\gu{કમ્પ્યૂટર}} — \emph{kampyūṭar} — \textquotedblleft{}computer\textquotedblright{}.

\begin{quote}
\textbf{\gu{અમારી} \gu{કચેરીમાં} \gu{નવું} \gu{કમ્પ્યૂટર} \gu{છે}.} — \emph{amārī kacerīmā̃ navũ kampyūṭar che} — \textquotedblleft{}There is a new computer in our office.\textquotedblright{}
\end{quote}

\begin{cousinweb}
\textbf{\gu{કમ્પ્યૂટર}} (\emph{kampyūṭar}) is a loan from English \emph{computer}.
\end{cousinweb}

\begin{grammarlens}[title={What you've built}]
One more word for everyday life at A2.
\end{grammarlens}

\subsection*{Your turn}
\begin{itemize}
\raggedright
  \item \emph{Say it:} \emph{kampyūṭar}
  \item \emph{Say it:} \emph{kampyūṭar}, once more
  \item \emph{Say it:} \emph{amārī kacerīmā̃ navũ kampyūṭar che} — There is a new computer in our office.
  \item \emph{Recall:} say the Gujarati for voice, then the Gujarati for to obtain, then say \emph{kampyūṭar} again
\end{itemize}

\begin{tcolorbox}[breakable,colback=teal!4,colframe=teal!35!black,title={Before you move on}]
What does \gu{કમ્પ્યૂટર} mean? (\textquotedblleft{}Computer\textquotedblright{}.) Say it once more.
\end{tcolorbox}
\section[\texorpdfstring{\gu{સ્ક્રીન} (skrīn)}{સ્ક્રીન (skrīn)}]{\gu{સ્ક્રીન} (skrīn) — screen}

\label{lesson:GU-C250-skrin}



\begin{quote}
Before the new one: say the Gujarati for to obtain, then the Gujarati for computer.
\end{quote}

\subsection*{You'll want to know: \gu{સ્ક્રીન}}
\textbf{\gu{સ્ક્રીન}} — \emph{skrīn} — \textquotedblleft{}screen\textquotedblright{}.

\begin{quote}
\textbf{\gu{કમ્પ્યૂટરની} \gu{સ્ક્રીન} \gu{મોટી} \gu{છે}.} — \emph{kampyūṭarnī skrīn moṭī che} — \textquotedblleft{}The computer's screen is big.\textquotedblright{}
\end{quote}

\begin{cousinweb}
\textbf{\gu{સ્ક્રીન}} (\emph{skrīn}) is a loan from English \emph{screen}, used for computers, phones and the cinema.
\end{cousinweb}

\begin{grammarlens}[title={What you've built}]
One more word for everyday life at A2.
\end{grammarlens}

\subsection*{Your turn}
\begin{itemize}
\raggedright
  \item \emph{Say it:} \emph{skrīn}
  \item \emph{Say it:} \emph{skrīn}, once more
  \item \emph{Say it:} \emph{kampyūṭarnī skrīn moṭī che} — The computer's screen is big.
  \item \emph{Recall:} say the Gujarati for to obtain, then the Gujarati for computer, then say \emph{skrīn} again
\end{itemize}

\begin{tcolorbox}[breakable,colback=teal!4,colframe=teal!35!black,title={Before you move on}]
What does \gu{સ્ક્રીન} mean? (\textquotedblleft{}Screen\textquotedblright{}.) Say it once more.
\end{tcolorbox}
\section[\texorpdfstring{\gu{દબાવવું} (dabāvvũ)}{દબાવવું (dabāvvũ)}]{\gu{દબાવવું} (dabāvvũ) — to press}

\label{lesson:GU-C250-dabavvu}



\begin{quote}
Before the new one: say the Gujarati for computer, then the Gujarati for screen.
\end{quote}

\subsection*{You'll want to know: \gu{દબાવવું}}
\textbf{\gu{દબાવવું}} — \emph{dabāvvũ} — \textquotedblleft{}to press\textquotedblright{}.

\begin{quote}
\textbf{\gu{તે} \gu{લાલ} \gu{બટન} \gu{દબાવે} \gu{છે}.} — \emph{te lāl baṭan dabāve che} — \textquotedblleft{}She presses the red button.\textquotedblright{}
\end{quote}

\begin{cousinweb}
\textbf{\gu{દબાવવું}} (\emph{dabāvvũ}) comes from \emph{*dabb}, \textquotedblleft{}press\textquotedblright{}, a reconstructed form that CDIAL, citing A. Master, traces to Dravidian, so it may not be inherited from Sanskrit at all; CDIAL glosses the Gujarati verb \textquotedblleft{}to press, intimidate\textquotedblright{} (CDIAL 6173).
\end{cousinweb}

\begin{grammarlens}[title={What you've built}]
One more word for everyday life at A2.
\end{grammarlens}

\subsection*{Your turn}
\begin{itemize}
\raggedright
  \item \emph{Say it:} \emph{dabāvvũ}
  \item \emph{Say it:} \emph{dabāvvũ}, once more
  \item \emph{Say it:} \emph{te lāl baṭan dabāve che} — She presses the red button.
  \item \emph{Recall:} say the Gujarati for computer, then the Gujarati for screen, then say \emph{dabāvvũ} again
\end{itemize}

\begin{tcolorbox}[breakable,colback=teal!4,colframe=teal!35!black,title={Before you move on}]
What does \gu{દબાવવું} mean? (\textquotedblleft{}To press\textquotedblright{}.) Say it once more.
\end{tcolorbox}
\section[\texorpdfstring{\gu{સાચવવું} (sācavvũ)}{સાચવવું (sācavvũ)}]{\gu{સાચવવું} (sācavvũ) — to keep safe, to save}

\label{lesson:GU-C250-sacavvu}



\begin{quote}
Before the new one: say the Gujarati for screen, then the Gujarati for to press.
\end{quote}

\subsection*{You'll want to know: \gu{સાચવવું}}
\textbf{\gu{સાચવવું}} — \emph{sācavvũ} — \textquotedblleft{}to keep safe, to save\textquotedblright{}.

\begin{quote}
\textbf{\gu{હું} \gu{મારું} \gu{કામ} \gu{કમ્પ્યૂટરમાં} \gu{સાચવું} \gu{છું}.} — \emph{hũ mārũ kām kampyūṭarmā̃ sācavũ chũ} — \textquotedblleft{}I save my work on the computer.\textquotedblright{}
\end{quote}

\begin{cousinweb}
\textbf{\gu{સાચવવું}} (\emph{sācavvũ}) is a tadbhava. CDIAL lists it, glossed \textquotedblleft{}to preserve\textquotedblright{}, under Sanskrit \emph{sájati}, \textquotedblleft{}is attached, clings\textquotedblright{}, as built on an analogical Middle Indo-Aryan stem \emph{*ṣacc-} (CDIAL 13085).
\end{cousinweb}

\begin{grammarlens}[title={What you've built}]
One more word for everyday life at A2.
\end{grammarlens}

\subsection*{Your turn}
\begin{itemize}
\raggedright
  \item \emph{Say it:} \emph{sācavvũ}
  \item \emph{Say it:} \emph{sācavvũ}, once more
  \item \emph{Say it:} \emph{hũ mārũ kām kampyūṭarmā̃ sācavũ chũ} — I save my work on the computer.
  \item \emph{Recall:} say the Gujarati for screen, then the Gujarati for to press, then say \emph{sācavvũ} again
\end{itemize}

\begin{tcolorbox}[breakable,colback=teal!4,colframe=teal!35!black,title={Before you move on}]
What does \gu{સાચવવું} mean? (\textquotedblleft{}To keep safe, to save\textquotedblright{}.) Say it once more.
\end{tcolorbox}
\section[\texorpdfstring{\gu{છાપવું} (chāpvũ)}{છાપવું (chāpvũ)}]{\gu{છાપવું} (chāpvũ) — to print}

\label{lesson:GU-C250-chapvu}



\begin{quote}
Before the new one: say the Gujarati for to press, then the Gujarati for to keep safe.
\end{quote}

\subsection*{You'll want to know: \gu{છાપવું}}
\textbf{\gu{છાપવું}} — \emph{chāpvũ} — \textquotedblleft{}to print\textquotedblright{}.

\begin{quote}
\textbf{\gu{તે} \gu{અરજી} \gu{છાપે} \gu{છે}.} — \emph{te arjī chāpe che} — \textquotedblleft{}He prints the application.\textquotedblright{}
\end{quote}

\begin{cousinweb}
\textbf{\gu{છાપવું}} (\emph{chāpvũ}) is a tadbhava. It comes from \emph{*chapp}, \textquotedblleft{}press, cover\textquotedblright{}, a reconstructed form; CDIAL glosses the Gujarati verb \textquotedblleft{}to print\textquotedblright{} (CDIAL 4994).
\end{cousinweb}

\begin{grammarlens}[title={What you've built}]
One more word for everyday life at A2.
\end{grammarlens}

\subsection*{Your turn}
\begin{itemize}
\raggedright
  \item \emph{Say it:} \emph{chāpvũ}
  \item \emph{Say it:} \emph{chāpvũ}, once more
  \item \emph{Say it:} \emph{te arjī chāpe che} — He prints the application.
  \item \emph{Recall:} say the Gujarati for to press, then the Gujarati for to keep safe, then say \emph{chāpvũ} again
\end{itemize}

\begin{tcolorbox}[breakable,colback=teal!4,colframe=teal!35!black,title={Before you move on}]
What does \gu{છાપવું} mean? (\textquotedblleft{}To print\textquotedblright{}.) Say it once more.
\end{tcolorbox}
